\documentclass[a4paper,11pt]{article}

\usepackage[spanish]{babel}
\usepackage[T1]{fontenc}
\usepackage[margin=2.2cm]{geometry}
\usepackage{amsmath,amssymb}
\usepackage[shortlabels]{enumitem}
\usepackage{xcolor}
\usepackage{parskip}

\newcommand{\respuesta}[1]{\par\textcolor{blue!60!black}{\textbf{Solución.}} #1}
\newcommand{\correcta}[1]{\textcolor{red!70!black}{\textbf{#1}}}

\title{Parcial 1 de Astronomía Práctica II}
\author{}
\date{Semestre 2026-2}

\begin{document}

\maketitle

\section*{1. Resolución angular}

Dos telescopios tienen apertura $D=100\,\mathrm{cm}=1\,\mathrm{m}$. El telescopio A observa en $\lambda_A=500\,\mathrm{nm}$ y el B en $\lambda_B=150\,\mathrm{nm}$.

\begin{enumerate}[a)]
    \item Calcule el límite de resolución de Rayleigh de ambos telescopios en segundos de arco.

    \respuesta{Se utiliza
    \[
        \theta_{\mathrm{lim}}=1{,}22\frac{\lambda}{D},
        \qquad
        1\,\mathrm{rad}=206265^{\prime\prime}.
    \]
    Por tanto,
    \[
        \theta_A=1{,}22\frac{500\times10^{-9}}{1}(206265)
        \approx0{,}126^{\prime\prime},
    \]
    \[
        \theta_B=1{,}22\frac{150\times10^{-9}}{1}(206265)
        \approx0{,}0377^{\prime\prime}.
    \]}

    \item ¿En qué factor es superior la resolución espacial del telescopio B respecto a la del A?

    \respuesta{
    \[
        \frac{\theta_A}{\theta_B}
        =\frac{\lambda_A}{\lambda_B}
        =\frac{500}{150}\approx3{,}33.
    \]
    El telescopio B puede distinguir detalles angulares aproximadamente $3{,}33$ veces menores.}

    \item ¿Qué impide aprovechar desde la superficie terrestre la ventaja de observar en el ultravioleta?

    \respuesta{La atmósfera absorbe fuertemente el UV, especialmente a $150\,\mathrm{nm}$. Además, la turbulencia atmosférica limita la resolución mediante el \textit{seeing}. Estas observaciones requieren normalmente telescopios espaciales.}
\end{enumerate}

\section*{2. Escala de placa y muestreo}

Un telescopio refractor tiene $f_{\mathrm{obj}}=1200\,\mathrm{mm}$ y $D=100\,\mathrm{mm}$. La CCD posee píxeles de $9\,\mu\mathrm{m}$.

\begin{enumerate}[a)]
    \item Calcule la escala de placa en segundos de arco por milímetro.

    \respuesta{
    \[
        P_s=\frac{206265}{f_{\mathrm{obj}}}
        =\frac{206265}{1200}
        \approx171{,}9^{\prime\prime}/\mathrm{mm}.
    \]}

    \item Calcule la escala en segundos de arco por píxel.

    \respuesta{Como $9\,\mu\mathrm{m}=0{,}009\,\mathrm{mm}$,
    \[
        P_{\mathrm{pix}}=P_s(0{,}009)
        \approx1{,}55^{\prime\prime}/\text{píxel}.
    \]}

    \item Si el límite de difracción es $1{,}4^{\prime\prime}$, ¿cuántos píxeles abarca un elemento mínimo de resolución?

    \respuesta{
    \[
        N_{\mathrm{pix}}
        =\frac{1{,}4^{\prime\prime}}{1{,}55^{\prime\prime}/\text{píxel}}
        \approx0{,}90\ \text{píxeles}.
    \]
    La imagen está submuestreada: el criterio de Nyquist requiere al menos dos píxeles por elemento de resolución.}
\end{enumerate}

\section*{3. Aumento mínimo y magnitud límite}

Un telescopio tiene $D=150\,\mathrm{mm}$, $f_{\mathrm{obj}}=1500\,\mathrm{mm}$ y se usa con una pupila ocular nocturna de $7\,\mathrm{mm}$.

\begin{enumerate}[a)]
    \item Calcule el aumento mínimo útil y la máxima distancia focal útil del ocular.

    \respuesta{
    \[
        M_{\min}=\frac{D}{d_{\mathrm{ojo}}}
        =\frac{150}{7}\approx21{,}4\times.
    \]
    Puesto que $M=f_{\mathrm{obj}}/f_{\mathrm{oc}}$,
    \[
        f_{\mathrm{oc,max}}
        =\frac{1500}{21{,}4}\approx70\,\mathrm{mm}.
    \]}

    \item Estime la magnitud límite estelar observable.

    \respuesta{Con $D_t=15\,\mathrm{cm}$ y una pupila de referencia de $0{,}6\,\mathrm{cm}$,
    \[
        m_{\lim}=6+5\log_{10}\left(\frac{15}{0{,}6}\right)
        \approx13{,}0.
    \]
    Es una estimación ideal para cielo oscuro y óptica sin pérdidas.}
\end{enumerate}

\section*{4. Consulta de Gaia}

Considere la consulta ADQL:

\begin{verbatim}
SELECT TOP 20 source_id, ra, dec, parallax, phot_g_mean_mag
FROM gaiadr3.gaia_source
WHERE parallax > 20 AND phot_g_mean_mag < 10
ORDER BY parallax DESC
\end{verbatim}

\begin{enumerate}[a)]
    \item Explique la consulta.

    \respuesta{Selecciona 20 fuentes de Gaia DR3 con paralaje mayor que $20\,\mathrm{mas}$ y magnitud media $G<10$. Devuelve su identificador, coordenadas, paralaje y magnitud, ordenadas desde el mayor paralaje; aproximadamente, aparecen primero las fuentes más cercanas.}

    \item Una estrella tiene paralaje $\pi=25\,\mathrm{mas}$. Calcule su distancia en parsecs y años luz.

    \respuesta{
    \[
        d(\mathrm{pc})=\frac{1000}{\pi(\mathrm{mas})}
        =\frac{1000}{25}=40\,\mathrm{pc}.
    \]
    Como $1\,\mathrm{pc}=3{,}2616$ años luz,
    \[
        d\approx40(3{,}2616)=130{,}5\ \text{años luz}.
    \]}
\end{enumerate}

\section*{5. Bases de datos astronómicas}

Indique la base de datos o archivo más apropiado para cada necesidad.

\begin{enumerate}[a)]
    \item Se desean obtener identificadores alternativos, coordenadas y bibliografía de Betelgeuse.\\
    \correcta{SIMBAD.}

    \item Se busca un catálogo publicado y se requieren filtros sobre sus columnas, por ejemplo Gaia DR3.\\
    \correcta{VizieR.}

    \item Se quieren descargar imágenes calibradas de Hubble o JWST de la galaxia M51.\\
    \correcta{MAST.}

    \item Se necesitan distancias, morfología y datos en múltiples longitudes de onda de M31.\\
    \correcta{NED.}
\end{enumerate}

\section*{6. Selección múltiple: IRAF}

\begin{enumerate}[a)]
    \item ¿Qué comando restablece los parámetros de una tarea a sus valores predeterminados?
    \begin{enumerate}[A.]
        \item \texttt{lpar}
        \item \texttt{epar}
        \item \correcta{\texttt{unlearn}}
        \item \texttt{flprcache}
    \end{enumerate}

    \item ¿Qué comando copia las columnas 180--220 y las filas 130--200 de \texttt{m51.fits} en \texttt{nucleo.fits}?
    \begin{enumerate}[A.]
        \item \texttt{imcopy m51.fits nucleo.fits[180:220,130:200]}
        \item \correcta{\texttt{imcopy m51.fits[180:220,130:200] nucleo.fits}}
        \item \texttt{display m51.fits[180:220,130:200] nucleo.fits}
        \item \texttt{imarith m51.fits[180:220,130:200] nucleo.fits}
    \end{enumerate}

    \item En el modo interactivo de \texttt{imexamine}, ¿qué tecla realiza fotometría de apertura sobre la fuente señalada?
    \begin{enumerate}[A.]
        \item \texttt{r}
        \item \correcta{\texttt{a}}
        \item \texttt{s}
        \item \texttt{q}
    \end{enumerate}
\end{enumerate}

\end{document}
